\documentclass[10pt]{article}
\usepackage{amsmath,amssymb,xcolor}
\usepackage[paperwidth=4.6in,paperheight=1.75in,margin=0.12in]{geometry}
\pagestyle{empty}
\begin{document}
\noindent Thin print fonts such as \emph{Computer Modern} were designed for ink on
paper. On screen, their hairline strokes can look faint, especially in
viewers that do not apply font smoothing. Stroking each glyph outline with
a very thin line restores some of the weight that ink spread used to add.
\[
  \mathcal{L}(\boldsymbol\phi) = \sum_{i=1}^{I} \bigl(\mathrm{f}[\mathbf{x}_i,\boldsymbol\phi] - y_i\bigr)^2,
  \qquad \hat{\boldsymbol\phi} = \operatorname*{argmin}_{\boldsymbol\phi}\, \mathcal{L}(\boldsymbol\phi).
\]
\textcolor{teal}{Coloured text keeps its colour.}
\end{document}
